\documentclass[12pt,a4paper]{article}
\usepackage[utf8]{inputenc}
\usepackage[T1]{fontenc}
\usepackage[brazil]{babel}
\usepackage{amsmath,amssymb,amsthm,mathtools}
\usepackage{graphicx}
\usepackage[margin=2.5cm,top=3cm,bottom=3cm]{geometry}
\usepackage{tikz}
\usetikzlibrary{arrows.meta,shapes.geometric,positioning,fit,backgrounds,calc,shadows.blur,decorations.pathmorphing,mindmap,trees,patterns,matrix}
\usepackage{pgfplots}
\pgfplotsset{compat=1.18}
\usepackage{fontawesome5}
\usepackage[ruled,vlined,linesnumbered]{algorithm2e}
\usepackage{listings}
\usepackage{xcolor}
\usepackage{booktabs}
\usepackage{multirow}
\usepackage{array}
\usepackage{tcolorbox}
\tcbuselibrary{breakable,skins,theorems}
\usepackage[hidelinks]{hyperref}
\usepackage{bookmark}
\setlength{\headheight}{14pt}
\usepackage{fancyhdr}
\usepackage{titlesec}
\usepackage{enumitem}
\usepackage{subcaption}
\usepackage{caption}
\definecolor{mineblue}{RGB}{0,90,180}
\definecolor{mineorange}{RGB}{230,115,0}
\definecolor{minegreen}{RGB}{34,139,34}
\definecolor{minered}{RGB}{180,20,40}
\definecolor{minegray}{RGB}{100,100,100}
\definecolor{codebg}{RGB}{250,250,250}
\newtcolorbox{definicaobox}[1]{title={\faIcon{book} #1},colback=blue!4!white,colframe=mineblue,fonttitle=\bfseries,breakable,enhanced}
\newtcolorbox{exemplobox}[1]{title={\faIcon{lightbulb} #1},colback=orange!4!white,colframe=mineorange,fonttitle=\bfseries,breakable,enhanced}
\newtcolorbox{dicabox}[1]{title={\faIcon{info-circle} #1},colback=green!4!white,colframe=minegreen,fonttitle=\bfseries,breakable,enhanced}
\newtcolorbox{atencaobox}[1]{title={\faIcon{exclamation-triangle} #1},colback=red!4!white,colframe=minered,fonttitle=\bfseries,breakable,enhanced}
\newtcolorbox{estadoartebox}[1]{title={\faIcon{rocket} #1},colback=purple!4!white,colframe=purple!60!black,fonttitle=\bfseries,breakable,enhanced}
\lstdefinestyle{mypython}{language=Python,backgroundcolor=\color{codebg},commentstyle=\color{minegray}\itshape,keywordstyle=\color{mineblue}\bfseries,stringstyle=\color{minegreen},basicstyle=\ttfamily\footnotesize,breaklines=true,frame=single,rulecolor=\color{minegray!40},numbers=left,numberstyle=\tiny\color{minegray},tabsize=4,showstringspaces=false}
\lstset{style=mypython}
\pagestyle{fancy}
\fancyhf{}
\fancyhead[L]{\small \textcolor{mineblue}{\textbf{Mineração de Dados}} | UEPA -- CCNT}
\fancyhead[R]{\small Módulo 6: Métodos Estatísticos e Clustering}
\fancyfoot[C]{\small \thepage}
\renewcommand{\headrulewidth}{0.4pt}
\titleformat{\section}{\large\bfseries\color{mineblue}}{\thesection}{0.8em}{}[\titlerule]
\titleformat{\subsection}{\normalsize\bfseries\color{mineorange}}{\thesubsection}{0.8em}{}
\titleformat{\subsubsection}{\small\bfseries\color{minegray}}{\thesubsubsection}{0.8em}{}

\begin{document}

\begin{center}
  {\LARGE\bfseries\color{mineblue} Módulo 6: Métodos Estatísticos}\\[0.4em]
  {\large\color{mineorange} e Mineração Não-Supervisionada}\\[0.6em]
  {\normalsize Universidade do Estado do Pará (UEPA) -- CCNT\\
  Curso: Engenharia de Software, 8$^{\circ}$ Semestre\\
  Disciplina: Mineração de Dados | Carga Horária: 80h}\\[0.3em]
  \rule{\linewidth}{0.8pt}
\end{center}

% ============================================================
\section{Objetivos de Aprendizagem}
% ============================================================

\begin{itemize}[leftmargin=2em]
  \item[\faIcon{check-circle}] \textbf{Aplicar} testes de hipótese estatísticos para comparar o desempenho de algoritmos de mineração de dados, interpretando p-valores e erros Tipo I e II.
  \item[\faIcon{check-circle}] \textbf{Calcular} e interpretar medidas de correlação e dependência (Pearson, Spearman, Informação Mútua) para análise exploratória de dados.
  \item[\faIcon{check-circle}] \textbf{Implementar} e comparar algoritmos de agrupamento (K-Means, Hierarchical, DBSCAN, GMM) para descoberta de estrutura em dados não rotulados.
  \item[\faIcon{check-circle}] \textbf{Selecionar} o algoritmo de clustering mais adequado conforme a forma geométrica dos grupos, presença de ruído e dimensionalidade dos dados.
  \item[\faIcon{check-circle}] \textbf{Avaliar} agrupamentos com métricas internas (Silhouette, Davies-Bouldin) e externas (ARI, NMI), compreendendo as limitações de cada abordagem.
  \item[\faIcon{check-circle}] \textbf{Minerar} regras de associação usando os algoritmos Apriori e FP-Growth, calculando suporte, confiança, lift e conviction.
  \item[\faIcon{check-circle}] \textbf{Detectar} anomalias usando métodos estatísticos, baseados em distância e em florestas de isolamento (Isolation Forest).
  \item[\faIcon{check-circle}] \textbf{Identificar} as tendências do estado da arte (2020--2024) em clustering profundo, HDBSCAN e detecção moderna de outliers.
\end{itemize}

% ============================================================
\section{Fundamentos Estatísticos para Mineração de Dados}
% ============================================================

\subsection{Testes de Hipótese}

Um \textbf{teste de hipótese} é um procedimento para decidir, com base em evidências amostrais, se uma afirmação sobre uma população deve ser aceita ou rejeitada.

\begin{definicaobox}{Componentes de um Teste de Hipótese}
\begin{itemize}
  \item \textbf{$H_0$ (hipótese nula)}: A posição padrão a ser testada (ex.\ ``algoritmo A e B têm a mesma acurácia'').
  \item \textbf{$H_1$ (hipótese alternativa)}: O que queremos provar (ex.\ ``algoritmo A é melhor que B'').
  \item \textbf{Nível de significância $\alpha$}: Probabilidade máxima do Erro Tipo I aceitável (tipicamente $\alpha = 0.05$).
  \item \textbf{p-valor}: Probabilidade de observar um resultado tão extremo quanto o observado, supondo $H_0$ verdadeira. Rejeita-se $H_0$ se $p < \alpha$.
\end{itemize}
\end{definicaobox}

\textbf{Tipos de Erro:}
\begin{itemize}
  \item \textbf{Erro Tipo I ($\alpha$)}: Rejeitar $H_0$ quando ela é verdadeira (\textit{falso positivo}).
  \item \textbf{Erro Tipo II ($\beta$)}: Não rejeitar $H_0$ quando ela é falsa (\textit{falso negativo}).
  \item \textbf{Poder do Teste ($1-\beta$)}: Probabilidade de detectar corretamente o efeito quando $H_1$ é verdadeira.
\end{itemize}

\textbf{Testes Relevantes para Mineração de Dados:}

\begin{itemize}
  \item \textbf{Teste $t$ pareado}: Compara médias de dois classificadores nos mesmos folds de validação cruzada. Assume normalidade das diferenças.
    \[
      t = \frac{\bar{d}}{s_d / \sqrt{n}}, \quad \bar{d} = \frac{1}{n}\sum_{i=1}^{n}(a_i - b_i)
    \]
  \item \textbf{ANOVA (Analysis of Variance)}: Compara três ou mais classificadores simultaneamente. Decompõe variância total em variância entre grupos e dentro dos grupos.
  \item \textbf{Teste de Wilcoxon Signed-Rank} (não-paramétrico): Alternativa ao teste $t$ pareado quando os dados não seguem distribuição normal. Ranqueia as diferenças absolutas e testa se as medianas são iguais.
  \item \textbf{McNemar's Test}: Compara dois classificadores pelas matrizes de confusão. Testa se o número de casos onde A acerta e B erra é diferente do número onde B acerta e A erra:
    \[
      \chi^2 = \frac{(b - c)^2}{b + c}
    \]
    onde $b$ = casos em que A acerta e B erra, $c$ = casos em que B acerta e A erra.
\end{itemize}

\begin{atencaobox}{Comparações Múltiplas}
Ao comparar $k$ classificadores com múltiplos testes de hipótese, a probabilidade de pelo menos um falso positivo é $1-(1-\alpha)^k$. Para 10 comparações com $\alpha=0.05$, isso é $\approx 40\%$! Use correções como \textbf{Bonferroni} ($\alpha' = \alpha/k$) ou o teste de \textbf{Friedman} + post-hoc de Nemenyi para comparações múltiplas.
\end{atencaobox}

\subsection{Correlação e Dependência}

\textbf{Coeficiente de Correlação de Pearson} ($r$): Mede relação linear entre duas variáveis contínuas:
\[
  r = \frac{\sum_{i=1}^{n}(x_i - \bar{x})(y_i - \bar{y})}{\sqrt{\sum_{i=1}^{n}(x_i - \bar{x})^2 \cdot \sum_{i=1}^{n}(y_i - \bar{y})^2}}
  \in [-1, 1]
\]
$r = 1$ (correlação positiva perfeita), $r = 0$ (sem correlação linear), $r = -1$ (negativa perfeita).

\textbf{Correlação de Spearman} ($\rho$): Aplica Pearson nos \textbf{postos} (\textit{ranks}) dos dados. Detecta monotonicidade, não apenas linearidade. Robusta a outliers e útil para dados ordinais.

\textbf{Tau de Kendall} ($\tau$): Baseado em concordância/discordância de pares. Mais robusto que Spearman para amostras pequenas.

\textbf{Informação Mútua} $I(X;Y)$: Mede dependência geral (não apenas linear) entre variáveis:
\[
  I(X; Y) = \sum_{x,y} p(x,y) \log \frac{p(x,y)}{p(x) p(y)}
\]
$I(X;Y) = 0$ iff $X$ e $Y$ são independentes. Muito usado em seleção de features (MRMR).

\subsection{Distribuições Relevantes}

\begin{table}[h!]
\centering
\caption{Distribuições de probabilidade frequentes em Mineração de Dados.}
\renewcommand{\arraystretch}{1.5}
\begin{tabular}{>{\bfseries}p{2.5cm} p{4cm} p{1.5cm} p{4.5cm}}
\toprule
Distribuição & FDP/FMP & Suporte & Uso em DM \\
\midrule
Normal & $\frac{1}{\sigma\sqrt{2\pi}} e^{-\frac{(x-\mu)^2}{2\sigma^2}}$ & $\mathbb{R}$ & Ruído, GMM, testes paramétricos \\
Bernoulli & $p^x (1-p)^{1-x}$ & $\{0,1\}$ & Classificação binária, Naive Bayes \\
Binomial & $\binom{n}{k} p^k (1-p)^{n-k}$ & $\{0,\ldots,n\}$ & Contagem de acertos, limiar de confiança \\
Poisson & $\frac{\lambda^k e^{-\lambda}}{k!}$ & $\mathbb{N}_0$ & Contagem de eventos raros, detecção de anomalias \\
Exponencial & $\lambda e^{-\lambda x}$ & $[0,+\infty)$ & Tempo entre eventos, detecção de outliers \\
$t$-Student & $\frac{\Gamma(\frac{\nu+1}{2})}{\sqrt{\nu\pi}\,\Gamma(\frac{\nu}{2})}\left(1+\frac{t^2}{\nu}\right)^{-\frac{\nu+1}{2}}$ & $\mathbb{R}$ & Testes com pequenas amostras \\
Qui-Quadrado & $\frac{x^{k/2-1}e^{-x/2}}{2^{k/2}\Gamma(k/2)}$ & $[0,+\infty)$ & Teste de independência (tabelas de contingência) \\
\bottomrule
\end{tabular}
\end{table}

% ============================================================
\section{Agrupamento (Clustering)}
% ============================================================

\subsection{Definição e Aplicações}

O \textbf{agrupamento} (clustering) é uma tarefa de aprendizado não-supervisionado que busca descobrir a estrutura natural dos dados, organizando exemplos similares em grupos (\textit{clusters}) sem o uso de rótulos pré-definidos.

\textbf{Objetivo formal}: Dado $\mathcal{X} = \{\mathbf{x}_1, \ldots, \mathbf{x}_N\} \subset \mathbb{R}^d$, encontrar uma partição $\mathcal{C} = \{C_1, \ldots, C_K\}$ tal que objetos dentro de cada $C_k$ sejam similares entre si e dissimilares aos de outros clusters.

\textbf{Aplicações:}
\begin{itemize}
  \item \textbf{Segmentação de clientes}: Grupos de comportamento de compra para marketing direcionado.
  \item \textbf{Detecção de anomalias}: Pontos que não pertencem a nenhum cluster são candidatos a outliers.
  \item \textbf{Sumarização de dados}: Representa o dataset pelos centróides dos clusters (compressão).
  \item \textbf{Modelagem de tópicos}: Agrupa documentos por tema (text mining).
  \item \textbf{Bioinformática}: Agrupamento de genes por padrão de expressão.
  \item \textbf{Compressão de imagens}: K-Means para quantização de cores.
\end{itemize}

\begin{dicabox}{Avaliação sem Ground Truth}
A falta de rótulos torna a avaliação de clustering mais desafiadora do que a classificação supervisionada. Métricas internas (Silhouette, Davies-Bouldin) medem a coesão e separação dos clusters, mas podem favorecer clustering simples. Sempre combine métricas numéricas com inspeção visual (PCA, t-SNE, UMAP) para validar resultados.
\end{dicabox}

\subsection{K-Means}

O K-Means (MacQueen, 1967; Lloyd, 1982) é o algoritmo de clustering mais popular. Minimiza a \textbf{inércia} (soma dos quadrados das distâncias intra-cluster):
\[
  J = \sum_{k=1}^{K} \sum_{\mathbf{x}_i \in C_k} \|\mathbf{x}_i - \boldsymbol{\mu}_k\|^2
\]
onde $\boldsymbol{\mu}_k = \frac{1}{|C_k|}\sum_{\mathbf{x}_i \in C_k} \mathbf{x}_i$ é o centróide do cluster $k$.

\begin{algorithm}[H]
\caption{K-Means (Algoritmo de Lloyd)}
\KwIn{Dados $\mathcal{X}$, número de clusters $K$, tolerância $\varepsilon$}
\KwOut{Partição $\{C_1,\ldots,C_K\}$ e centróides $\{\boldsymbol{\mu}_1,\ldots,\boldsymbol{\mu}_K\}$}
Inicializar centróides $\boldsymbol{\mu}_1, \ldots, \boldsymbol{\mu}_K$ aleatoriamente (ou com K-Means++)\;
\Repeat{centróides não mudam (ou $\Delta < \varepsilon$)}{
  \tcp{Passo de Atribuição (E-step)}
  \ForEach{$\mathbf{x}_i \in \mathcal{X}$}{
    $c_i \leftarrow \arg\min_{k \in \{1,\ldots,K\}} \|\mathbf{x}_i - \boldsymbol{\mu}_k\|^2$\;
  }
  \tcp{Passo de Atualização (M-step)}
  \ForEach{$k \in \{1,\ldots,K\}$}{
    $\boldsymbol{\mu}_k \leftarrow \frac{1}{|C_k|} \sum_{\mathbf{x}_i : c_i = k} \mathbf{x}_i$\;
  }
}
\end{algorithm}

\textbf{Visualização das etapas do K-Means:}

\begin{center}
\begin{tikzpicture}[scale=0.75]
% ---- Painel 1: Inicialização ----
\begin{scope}[xshift=0cm]
\node[above,font=\small\bfseries,mineblue] at (2.5,4.2) {1. Inicialização};
% Cluster azul (canto superior esquerdo)
\foreach \x/\y in {0.8/3.2, 1.3/2.6, 0.6/3.7, 1.5/3.0, 1.0/2.8}{
  \filldraw[mineblue!50] (\x,\y) circle(0.09);
}
% Cluster laranja (canto superior direito)
\foreach \x/\y in {3.2/3.4, 3.8/2.8, 3.5/3.8, 4.0/3.2, 3.6/2.6}{
  \filldraw[mineorange!50] (\x,\y) circle(0.09);
}
% Cluster verde (baixo centro)
\foreach \x/\y in {2.0/0.8, 2.5/1.3, 1.8/1.5, 2.3/0.6, 2.8/1.0}{
  \filldraw[minegreen!50] (\x,\y) circle(0.09);
}
% Centróides aleatórios (vermelho)
\filldraw[minered] (2.5,2.5) node[above,font=\tiny]{$\mu_1$} circle(0.15);
\filldraw[minered] (1.8,2.0) node[above,font=\tiny]{$\mu_2$} circle(0.15);
\filldraw[minered] (3.0,1.5) node[above,font=\tiny]{$\mu_3$} circle(0.15);
\draw[minegray,dashed] (-0.2,0.0) rectangle (5.0,4.5);
\end{scope}

% ---- Painel 2: Primeira Atribuição ----
\begin{scope}[xshift=5.5cm]
\node[above,font=\small\bfseries,mineblue] at (2.5,4.2) {2. Primeira Atribuição};
\foreach \x/\y in {0.8/3.2, 1.3/2.6, 0.6/3.7, 1.5/3.0, 1.0/2.8}{
  \filldraw[mineblue!70] (\x,\y) circle(0.09);
}
\foreach \x/\y in {3.2/3.4, 3.8/2.8, 3.5/3.8, 4.0/3.2, 3.6/2.6}{
  \filldraw[mineorange!70] (\x,\y) circle(0.09);
}
\foreach \x/\y in {2.0/0.8, 2.5/1.3, 1.8/1.5, 2.3/0.6, 2.8/1.0}{
  \filldraw[minegreen!70] (\x,\y) circle(0.09);
}
\filldraw[mineblue!80] (1.04,3.06) node[above,font=\tiny,mineblue]{$\mu_1$} circle(0.15);
\filldraw[mineorange!80] (3.62,3.16) node[above,font=\tiny,mineorange]{$\mu_2$} circle(0.15);
\filldraw[minegreen!80] (2.28,1.04) node[right,font=\tiny,minegreen]{$\mu_3$} circle(0.15);
\draw[minegray,dashed] (-0.2,0.0) rectangle (5.0,4.5);
\end{scope}

% ---- Painel 3: Convergência ----
\begin{scope}[xshift=11cm]
\node[above,font=\small\bfseries,minegreen] at (2.5,4.2) {3. Convergência};
\foreach \x/\y in {0.8/3.2, 1.3/2.6, 0.6/3.7, 1.5/3.0, 1.0/2.8}{
  \filldraw[mineblue] (\x,\y) circle(0.09);
}
\foreach \x/\y in {3.2/3.4, 3.8/2.8, 3.5/3.8, 4.0/3.2, 3.6/2.6}{
  \filldraw[mineorange] (\x,\y) circle(0.09);
}
\foreach \x/\y in {2.0/0.8, 2.5/1.3, 1.8/1.5, 2.3/0.6, 2.8/1.0}{
  \filldraw[minegreen] (\x,\y) circle(0.09);
}
\filldraw[mineblue] (1.04,3.06) node[star,star points=5,fill=mineblue,draw=mineblue,minimum size=0.3cm,font=\tiny,above]{$\mu_1$} circle(0.15);
\filldraw[mineorange] (3.62,3.16) node[above,font=\tiny,mineorange]{$\mu_2$} circle(0.15);
\filldraw[minegreen] (2.28,1.04) node[right,font=\tiny,minegreen]{$\mu_3$} circle(0.15);
\draw[minegreen,dashed,thick] (-0.2,0.0) rectangle (5.0,4.5);
\node[below,font=\tiny\itshape,minegray] at (2.5,-0.1) {Centróides estabilizados};
\end{scope}
\end{tikzpicture}
\end{center}

\textbf{Escolha de K:}
\begin{itemize}
  \item \textbf{Método do Cotovelo (Elbow)}: Plota inércia vs $K$; o ``cotovelo'' indica o $K$ ótimo.
  \item \textbf{Análise de Silhouette}: Para cada $K$, calcula o coeficiente médio; escolhe $K$ com maior silhouette.
  \item \textbf{Gap Statistic} (Tibshirani et al., 2001): Compara inércia com a esperada sob distribuição de referência.
\end{itemize}

\textbf{Variantes do K-Means:}
\begin{itemize}
  \item \textbf{K-Means++} (Arthur \& Vassilvitskii, 2007): Inicialização melhorada -- cada centróide é escolhido com probabilidade proporcional à distância ao centróide mais próximo já escolhido. Garante aproximação $O(\log K)$ do ótimo.
  \item \textbf{K-Medoids (PAM)}: Usa exemplos reais do dataset como centróides (medóides). Mais robusto a outliers.
  \item \textbf{Mini-Batch K-Means}: Processa mini-lotes aleatórios; escala para datasets de milhões de exemplos.
\end{itemize}

\subsection{Clustering Hierárquico}

O clustering hierárquico cria uma \textbf{hierarquia} de agrupamentos representada em um \textbf{dendrograma}, sem necessidade de especificar $K$ antecipadamente.

\textbf{Agglomerativo (Bottom-up)}: Começa com $N$ clusters (um por ponto) e mescla iterativamente os dois mais próximos.

\textbf{Divisivo (Top-down)}: Começa com um único cluster e divide recursivamente.

\textbf{Critérios de Ligação (Linkage):}
\begin{align}
  d_{\text{single}}(A,B) &= \min_{a \in A, b \in B} d(a,b) && \text{(single linkage -- sensível a outliers)}\\
  d_{\text{complete}}(A,B) &= \max_{a \in A, b \in B} d(a,b) && \text{(complete linkage -- clusters compactos)}\\
  d_{\text{average}}(A,B) &= \frac{1}{|A||B|} \sum_{a \in A} \sum_{b \in B} d(a,b) && \text{(UPGMA -- balanço)}\\
  d_{\text{ward}}(A,B) &= \Delta J \text{ ao mesclar } A \text{ e } B && \text{(Ward -- minimiza variância intra-cluster)}
\end{align}

\textbf{Dendrograma (exemplo esquemático):}

\begin{center}
\begin{tikzpicture}[scale=0.85,
  level 1/.style={sibling distance=4cm, level distance=2cm},
  level 2/.style={sibling distance=2cm, level distance=1.5cm},
  level 3/.style={sibling distance=1cm, level distance=1.2cm},
  edge from parent/.style={draw,mineblue,thick}]
\node[draw=mineblue,fill=mineblue!10,rounded corners,font=\small\bfseries] {Raiz}
  child { node[draw=mineorange,fill=mineorange!10,rounded corners,font=\small] {Cluster A}
    child { node[draw=minegreen,fill=minegreen!10,rounded corners,font=\small] {A1}
      child { node[font=\footnotesize] {$p_1$} }
      child { node[font=\footnotesize] {$p_2$} }
    }
    child { node[draw=minegreen,fill=minegreen!10,rounded corners,font=\small] {A2}
      child { node[font=\footnotesize] {$p_3$} }
      child { node[font=\footnotesize] {$p_4$} }
    }
  }
  child { node[draw=mineorange,fill=mineorange!10,rounded corners,font=\small] {Cluster B}
    child { node[font=\footnotesize] {$p_5$} }
    child { node[draw=minegreen,fill=minegreen!10,rounded corners,font=\small] {B2}
      child { node[font=\footnotesize] {$p_6$} }
      child { node[font=\footnotesize] {$p_7$} }
    }
  };
\end{tikzpicture}
\end{center}

A \textbf{correlação cofenética} mede o quão fiel o dendrograma é às distâncias originais -- valores $> 0.75$ indicam boa representação hierárquica.

\subsection{DBSCAN}

O \textbf{DBSCAN} (Density-Based Spatial Clustering of Applications with Noise; Ester et al., 1996) é um algoritmo baseado em densidade que encontra clusters de \textbf{forma arbitrária} e detecta \textbf{ruído/outliers} automaticamente.

\textbf{Conceitos:} Dado $\varepsilon > 0$ e $\text{MinPts} \in \mathbb{N}$:
\begin{itemize}
  \item \textbf{Ponto Núcleo}: $|N_\varepsilon(\mathbf{x})| \geq \text{MinPts}$, onde $N_\varepsilon(\mathbf{x}) = \{\mathbf{y} : d(\mathbf{x},\mathbf{y}) \leq \varepsilon\}$.
  \item \textbf{Ponto de Fronteira}: Está na $\varepsilon$-vizinhança de um núcleo, mas não é núcleo.
  \item \textbf{Ponto de Ruído}: Não é núcleo nem fronteira.
\end{itemize}

\begin{center}
\begin{tikzpicture}[scale=1.0]
% Cluster 1 (denso, canto inferior esquerdo)
\foreach \x/\y in {1.5/1.0, 2.0/1.5, 1.0/1.5, 2.5/1.0, 1.8/2.0, 1.2/0.8, 2.2/0.8, 2.8/1.5}{
  \filldraw[mineblue!50] (\x,\y) circle(0.08);
}
% Mostrar epsilon para ponto núcleo
\draw[mineblue,dashed,thick] (1.8,1.4) circle(0.8);
\filldraw[mineblue] (1.8,1.4) circle(0.1);
\node[font=\tiny,above=0.85cm,mineblue] at (1.8,1.4) {ponto núcleo ($\varepsilon$)};

% Cluster 2 (topo direita)
\foreach \x/\y in {4.0/3.0, 4.5/3.5, 5.0/3.0, 4.8/2.5, 3.8/2.8, 5.2/3.5}{
  \filldraw[mineorange!60] (\x,\y) circle(0.08);
}

% Ponto de fronteira
\filldraw[minegreen,thick] (3.2/1,2.0/1) circle(0.1) node[right,font=\tiny,minegreen] {\;fronteira};

% Pontos de ruído
\filldraw[minered] (6.5/1,1.0/1) circle(0.1) node[right,font=\tiny,minered] {\;ruído};
\filldraw[minered] (0.3,3.5) circle(0.1) node[right,font=\tiny,minered] {\;ruído};

\node[below,font=\small\itshape,minegray] at (3.0,0.1) {DBSCAN: pontos núcleo (azul/laranja), fronteira (verde) e ruído (vermelho)};

% Legenda
\draw[mineblue,dashed,thick] (5.8,0.5) circle(0.2);
\node[right,font=\tiny] at (6.0,0.5) {$\varepsilon$-vizinhança};
\end{tikzpicture}
\end{center}

\begin{algorithm}[H]
\caption{DBSCAN}
\KwIn{Dados $\mathcal{X}$, raio $\varepsilon$, densidade mínima MinPts}
\KwOut{Rótulos de cluster $C$ para cada ponto (incluindo ruído $= -1$)}
$C \leftarrow 0$; \quad marcar todos os pontos como não visitados\;
\ForEach{ponto não visitado $\mathbf{p} \in \mathcal{X}$}{
  Marcar $\mathbf{p}$ como visitado\;
  $N \leftarrow N_\varepsilon(\mathbf{p})$ \tcp{vizinhança $\varepsilon$}
  \If{$|N| < \text{MinPts}$}{
    Marcar $\mathbf{p}$ como ruído\;
  }\Else{
    $C \leftarrow C + 1$\;
    ExpandCluster($\mathbf{p}$, $N$, $C$, $\varepsilon$, MinPts)\;
  }
}

\SetKwFunction{FExpand}{ExpandCluster}
\SetKwProg{Fn}{Função}{:}{}
\Fn{\FExpand{$\mathbf{p}$, $N$, $C$, $\varepsilon$, MinPts}}{
  Adicionar $\mathbf{p}$ ao cluster $C$\;
  \ForEach{$\mathbf{q} \in N$}{
    \If{$\mathbf{q}$ não visitado}{
      Marcar $\mathbf{q}$ como visitado\;
      $N' \leftarrow N_\varepsilon(\mathbf{q})$\;
      \If{$|N'| \geq \text{MinPts}$}{$N \leftarrow N \cup N'$\;}
    }
    \If{$\mathbf{q}$ ainda não atribuído a cluster}{Adicionar $\mathbf{q}$ ao cluster $C$\;}
  }
}
\end{algorithm}

\textbf{Parâmetros:} $\varepsilon$ pode ser estimado pelo gráfico da $k$-NN distance (joelho da curva para $k = \text{MinPts}$). Regra prática: $\text{MinPts} \geq 2d$ onde $d$ é a dimensionalidade.

\subsection{HDBSCAN}

O \textbf{HDBSCAN} (Hierarchical DBSCAN; McInnes et al., 2017) estende o DBSCAN com uma abordagem hierárquica que \textbf{varia o limiar de densidade} automaticamente.

\textbf{Vantagens sobre DBSCAN:}
\begin{itemize}
  \item Detecta clusters com \textbf{densidades variáveis} (problema do DBSCAN: apenas $\varepsilon$ global).
  \item Mais robusto na seleção de parâmetros: apenas \texttt{min\_cluster\_size} é crítico.
  \item Constrói uma árvore de condensação de clusters e extrai os mais persistentes.
  \item Retorna probabilidade de pertencimento ao cluster (métrica de membrana).
\end{itemize}

\subsection{Modelos de Mistura Gaussiana (GMM)}

O GMM modela os dados como uma \textbf{mistura} de $K$ distribuições Gaussianas:
\[
  p(\mathbf{x}) = \sum_{k=1}^{K} \pi_k \, \mathcal{N}(\mathbf{x} \mid \boldsymbol{\mu}_k, \boldsymbol{\Sigma}_k)
\]
onde $\pi_k$ são os pesos da mistura ($\sum_k \pi_k = 1$, $\pi_k > 0$).

O GMM é treinado pelo \textbf{algoritmo EM} (Dempster et al., 1977):

\textbf{E-step} -- Calcular responsabilidades:
\[
  r_{ik} = \frac{\pi_k \, \mathcal{N}(\mathbf{x}_i \mid \boldsymbol{\mu}_k, \boldsymbol{\Sigma}_k)}{\sum_{j=1}^{K} \pi_j \, \mathcal{N}(\mathbf{x}_i \mid \boldsymbol{\mu}_j, \boldsymbol{\Sigma}_j)}
\]

\textbf{M-step} -- Atualizar parâmetros:
\[
  \pi_k = \frac{1}{N}\sum_i r_{ik}, \quad
  \boldsymbol{\mu}_k = \frac{\sum_i r_{ik} \mathbf{x}_i}{\sum_i r_{ik}}, \quad
  \boldsymbol{\Sigma}_k = \frac{\sum_i r_{ik}(\mathbf{x}_i - \boldsymbol{\mu}_k)(\mathbf{x}_i - \boldsymbol{\mu}_k)^T}{\sum_i r_{ik}}
\]

\textbf{Seleção do número de componentes} $K$: Use \textbf{BIC} (Bayesian Information Criterion) ou \textbf{AIC}: $\text{BIC} = -2\log\hat{L} + p \log N$, onde $p$ é o número de parâmetros. Minimize o BIC.

\subsection{Spectral Clustering}

O Spectral Clustering (Ng et al., 2002) baseia-se na teoria de grafos espectrais para capturar estruturas não-convexas que K-Means falha.

\textbf{Algoritmo:}
\begin{enumerate}
  \item Construir matriz de similaridade $\mathbf{W}$ (ex.\ kernel RBF: $w_{ij} = e^{-\|x_i-x_j\|^2/2\sigma^2}$).
  \item Calcular Laplaciano normalizado: $\mathbf{L} = \mathbf{D}^{-1/2}(\mathbf{D}-\mathbf{W})\mathbf{D}^{-1/2}$, onde $D_{ii} = \sum_j w_{ij}$.
  \item Extrair os $K$ menores autovetores de $\mathbf{L}$; formar matriz $\mathbf{U}$.
  \item Aplicar K-Means no espaço dos autovetores.
\end{enumerate}

\subsection{Avaliação de Agrupamentos}

\subsubsection{Métricas Internas (sem rótulos verdadeiros)}

\textbf{Coeficiente de Silhouette} (Rousseeuw, 1987): Para cada ponto $i$:
\[
  s(i) = \frac{b(i) - a(i)}{\max\{a(i), b(i)\}}
\]
onde $a(i) = $ distância média intra-cluster, $b(i) = $ distância média ao cluster vizinho mais próximo.
$s(i) \in [-1, 1]$: valores próximos a 1 indicam boa alocação.

\textbf{Índice Davies-Bouldin} (Davies \& Bouldin, 1979):
\[
  \text{DB} = \frac{1}{K} \sum_{k=1}^{K} \max_{k \neq k'} \frac{s_k + s_{k'}}{\|c_k - c_{k'}\|}
\]
onde $s_k$ é a distância média intra-cluster $k$. Minimizar DB.

\textbf{Índice Calinski-Harabasz}: Razão entre dispersão inter-cluster e intra-cluster. Maximizar.

\subsubsection{Métricas Externas (com rótulos verdadeiros)}

\textbf{Adjusted Rand Index (ARI)}: Versão ajustada para chance do Rand Index. $\text{ARI} \in [-1, 1]$; 1 = agrupamento perfeito; 0 = aleatório.

\textbf{Normalized Mutual Information (NMI)}:
\[
  \text{NMI}(U, V) = \frac{2 \cdot I(U; V)}{H(U) + H(V)} \in [0, 1]
\]

\textbf{V-Measure}: Média harmônica de homogeneidade (cada cluster contém apenas uma classe) e completude (todos os membros de uma classe estão no mesmo cluster).

% ============================================================
\section{Regras de Associação}
% ============================================================

\subsection{Conceitos Fundamentais}

Dada uma base de transações $\mathcal{T} = \{T_1, \ldots, T_N\}$ sobre um conjunto de itens $\mathcal{I}$, o objetivo é encontrar regras $A \Rightarrow B$ ($A, B \subseteq \mathcal{I}$, $A \cap B = \emptyset$) com alta frequência e confiabilidade.

\textbf{Medidas de interesse:}
\begin{align}
  \text{sup}(A) &= \frac{|\{T \in \mathcal{T} : A \subseteq T\}|}{|\mathcal{T}|} && \text{(frequência do itemset)}\\
  \text{sup}(A \Rightarrow B) &= \text{sup}(A \cup B) && \text{(frequência da regra)}\\
  \text{conf}(A \Rightarrow B) &= \frac{\text{sup}(A \cup B)}{\text{sup}(A)} && \text{(confiança)}\\
  \text{lift}(A \Rightarrow B) &= \frac{\text{conf}(A \Rightarrow B)}{\text{sup}(B)} && \text{(importância relativa)}\\
  \text{conv}(A \Rightarrow B) &= \frac{1 - \text{sup}(B)}{1 - \text{conf}(A \Rightarrow B)} && \text{(convicção)}
\end{align}

\textbf{Interpretação:}
\begin{itemize}
  \item \textbf{Lift $> 1$}: $A$ e $B$ são positivamente correlacionados; a regra é mais frequente que o esperado por acaso.
  \item \textbf{Lift $= 1$}: $A$ e $B$ são independentes.
  \item \textbf{Lift $< 1$}: $A$ e $B$ são negativamente correlacionados (regra negativa).
\end{itemize}

\begin{exemplobox}{Exemplo: Cesta de Mercado}
Considere 5 transações: $T_1=\{leite, pão\}$, $T_2=\{leite, manteiga, pão\}$, $T_3=\{manteiga, queijo\}$, $T_4=\{leite, manteiga\}$, $T_5=\{leite, pão, queijo\}$.

Para a regra $\{leite\} \Rightarrow \{pão\}$:
\[
  \text{sup}(\{leite, pão\}) = 3/5 = 0.6, \quad
  \text{sup}(\{leite\}) = 4/5 = 0.8, \quad
  \text{conf} = 0.6/0.8 = 0.75
\]
\[
  \text{sup}(\{pão\}) = 3/5 = 0.6, \quad
  \text{lift} = 0.75/0.6 = 1.25 \quad \text{(correlação positiva)}
\]
\end{exemplobox}

\subsection{Algoritmo Apriori (Agrawal \& Srikant, 1994)}

O Apriori usa a propriedade \textbf{anti-monótona}: se um itemset é infrequente, todos os seus superconjuntos também são. Isso permite poda da busca.

\begin{algorithm}[H]
\caption{Apriori}
\KwIn{Transações $\mathcal{T}$, suporte mínimo $\text{minsup}$, confiança mínima $\text{minconf}$}
\KwOut{Conjunto de regras de associação frequentes}
$L_1 \leftarrow$ \{itemsets frequentes de tamanho 1\}\;
$k \leftarrow 1$\;
\While{$L_k \neq \emptyset$}{
  $C_{k+1} \leftarrow$ AprioriGen($L_k$) \tcp{gera candidatos de tamanho $k+1$}
  \ForEach{transação $T \in \mathcal{T}$}{
    $D_T \leftarrow$ \{candidatos em $C_{k+1}$ contidos em $T$\}\;
    Incrementar contagem de cada itemset em $D_T$\;
  }
  $L_{k+1} \leftarrow$ \{$c \in C_{k+1}$ : sup$(c) \geq \text{minsup}$\}\;
  $k \leftarrow k + 1$\;
}
$\mathcal{F} \leftarrow \bigcup_k L_k$ \tcp{todos os itemsets frequentes}
\ForEach{$f \in \mathcal{F}$, $A \subset f$, $B = f \setminus A$, $A \neq \emptyset$}{
  \If{conf$(A \Rightarrow B) \geq \text{minconf}$}{
    Adicionar $A \Rightarrow B$ às regras\;
  }
}
\end{algorithm}

\subsection{Algoritmo FP-Growth (Han et al., 2000)}

O FP-Growth elimina a geração de candidatos construindo uma estrutura compacta (\textbf{FP-Tree}) que representa a base de transações.

\textbf{Vantagens sobre Apriori:}
\begin{itemize}
  \item Apenas \textbf{dois varrimentos} da base de dados (independente do número de itemsets).
  \item Sem geração explícita de candidatos.
  \item A FP-Tree é tipicamente muito menor que a base original (compartilha prefixos).
\end{itemize}

\textbf{Funcionamento:} (1) Varre a base, conta frequências de itens, ordena por frequência decrescente. (2) Varre novamente, inserindo transações na FP-Tree. (3) Extrai itemsets frequentes recursivamente usando bases de padrões condicionais.

\textbf{ECLAT (Equivalence CLAss Transformation)}: Usa formato vertical dos dados (lista de transações por item -- TidList). Interseção de TidLists para computar suporte; eficiente para muitas transações.

\subsection{Regras de Associação para Cenários Específicos}

\begin{itemize}
  \item \textbf{Regras Quantitativas}: Itens com atributos numéricos; discretização ou particionamento fuzzy.
  \item \textbf{Regras Temporais / Mineração de Sequências}: Regras com ordenação temporal; algoritmo GSP (Generalized Sequential Patterns) e PrefixSpan.
  \item \textbf{Regras Multi-nível}: Taxonomias de itens (ex.\ leite $\subset$ laticínios $\subset$ alimentos); suporte mínimo ajustado por nível.
  \item \textbf{Regras Negativas}: $A \Rightarrow \neg B$ (comprar A implica não comprar B); úteis para descoberta de anti-correlações.
\end{itemize}

% ============================================================
\section{Detecção de Anomalias}
% ============================================================

\begin{definicaobox}{Tipos de Anomalias}
\begin{itemize}
  \item \textbf{Ponto}: Uma instância individual diverge significativamente das demais (ex.\ transação fraudulenta).
  \item \textbf{Contextual}: Uma instância é anômala em um contexto específico (ex.\ temperatura de $35^\circ$C em dezembro no RS).
  \item \textbf{Coletiva}: Um subconjunto de instâncias é anômalo em conjunto, mas cada instância individualmente parece normal.
\end{itemize}
\end{definicaobox}

\textbf{Métodos de Detecção de Anomalias:}

\begin{itemize}
  \item \textbf{Estatísticos}: Assume distribuição Gaussiana; Z-score ($|z| > 3$); Grubbs Test; Boxplot ($> 1.5 \times \text{IQR}$).
  \item \textbf{Baseados em Distância}: K-NN Outlier Score: anomalias têm vizinhos distantes.
  \item \textbf{Local Outlier Factor (LOF)} (Breunig et al., 2000): Compara a densidade de um ponto com a de seus vizinhos. $\text{LOF}(p) \gg 1$ indica anomalia.
  \item \textbf{Isolation Forest} (Liu et al., 2008): Constrói árvores de isolamento aleatórias. Anomalias são isoladas em menos partições (caminhos curtos na árvore). Score de anomalia: $s(x,n) = 2^{-E[h(x)]/c(n)}$, onde $h(x)$ é o comprimento do caminho e $c(n)$ é o comprimento médio esperado.
  \item \textbf{Autoencoder}: Treina a rede neural para reconstruir dados normais. Alta \textbf{erro de reconstrução} $\|x - \hat{x}\|^2$ indica anomalia.
  \item \textbf{VAE (Variational Autoencoder)}: Score de anomalia baseado no ELBO (Evidence Lower BOund).
\end{itemize}

% ============================================================
\section{Estado da Arte (2020--2024)}
% ============================================================

\begin{estadoartebox}{Deep Clustering}
Integração de redes neurais profundas com objetivos de clustering:
\begin{itemize}
  \item \textbf{DEC -- Deep Embedded Clustering} (Xie et al., 2016): Autoencoder treinado conjuntamente com uma função de perda de clustering; aprende representações otimizadas para agrupamento.
  \item \textbf{SCAN} (Van Gansbeke et al., 2020): Self-supervised clustering; usa similaridade de vizinhos para treinar sem rótulos; state-of-the-art em clustering de imagens.
  \item \textbf{DINO} (Caron et al., 2021): Recursos do Vision Transformer (ViT) pré-treinados auto-supervisionados revelam estrutura de clustering sem fine-tuning; propriedade emergente dos Transformers.
\end{itemize}
\end{estadoartebox}

\begin{estadoartebox}{HDBSCAN e Clustering para Dados de Alta Dimensão}
\begin{itemize}
  \item \textbf{HDBSCAN} tornou-se o padrão para clustering robusto em Python (scikit-learn $\geq 1.3$). Suporta dados de alta dimensão com a métrica de alcançabilidade mútua.
  \item \textbf{Pipeline UMAP + HDBSCAN}: UMAP reduz dimensionalidade preservando estrutura local, HDBSCAN agrupa no espaço reduzido. Muito eficaz para NLP (embeddings de documentos), genômica single-cell (scRNA-seq).
  \item \textbf{BERTopic} (Grootendorst, 2022): Modela tópicos combinando BERT embeddings + UMAP + HDBSCAN + TF-IDF por cluster. Supera LDA em qualidade de tópicos.
\end{itemize}
\end{estadoartebox}

\begin{estadoartebox}{Detecção de Outliers Moderna}
\begin{itemize}
  \item \textbf{PyOD} (Zhao et al., 2019): Biblioteca Python com 40+ algoritmos de detecção de outliers unificados em uma API sklearn-compatível.
  \item \textbf{COPOD} (Li et al., 2020): Copula-Based Outlier Detection; sem hiperparâmetros; muito eficiente; baseado em cópulas empíricas.
  \item \textbf{ECOD} (Li et al., 2022): Empirical Cumulative distribution-based Outlier Detection; computacionalmente eficiente; interpretável; sem hiperparâmetros.
  \item \textbf{Detecção de Anomalias em Grafos}: GraphSAGE + Isolation Forest para detecção de fraudes em redes sociais; muito relevante para Social CRM.
\end{itemize}
\end{estadoartebox}

% ============================================================
\section{Exercícios}
% ============================================================

\begin{enumerate}
  \item \textbf{Silhouette Manual}: Dados os pontos $p_1=(1,1), p_2=(2,1), p_3=(5,5), p_4=(6,5)$ e os clusters $C_1 = \{p_1, p_2\}$, $C_2 = \{p_3, p_4\}$, calcule o coeficiente de silhouette para cada ponto e a silhouette média. O clustering é bom?

  \item \textbf{K-Means com Diferentes K}: Aplique K-Means com $K = 2, 3, 4, 5$ ao dataset Iris (4 features). Para cada $K$, calcule a inércia, Silhouette e Davies-Bouldin. Plote as curvas e determine o $K$ ótimo. O $K$ ótimo coincide com o número real de classes (3)?

  \item \textbf{Comparação de Linkages}: Aplique clustering hierárquico com single, complete, average e Ward linkage ao dataset Wine. Para cada método, corte o dendrograma em $K=3$ e calcule ARI e NMI em relação às classes verdadeiras. Qual linkage tem melhor desempenho? Por quê?

  \item \textbf{Apriori Manual}: Dada a base de transações: $T_1=\{A,B,C\}$, $T_2=\{A,B\}$, $T_3=\{B,C,D\}$, $T_4=\{A,C\}$, $T_5=\{A,B,D\}$, $T_6=\{B,C\}$. Com suporte mínimo $= 0.5$ e confiança mínima $= 0.7$, execute o Apriori: liste os itemsets frequentes $L_1, L_2, L_3$ e as regras válidas com suporte, confiança e lift.

  \item \textbf{DBSCAN Parameter Selection}: Aplique DBSCAN ao dataset make\_moons de scikit-learn. Varie $\varepsilon \in \{0.1, 0.2, 0.3, 0.5\}$ e MinPts $\in \{3, 5, 10\}$. Para cada combinação, reporte número de clusters, número de pontos de ruído e ARI. Qual combinação é melhor? Compare com K-Means ($K=2$).

  \item \textbf{Comparação de Métodos de Anomalia}: Usando o dataset Pima Indians Diabetes ou Credit Card Fraud, compare Isolation Forest, LOF e Z-Score para detecção de anomalias. Avalie com AUC-ROC e Precisão@K (K = 1\% dos dados). Qual método tem melhor desempenho?

  \item \textbf{Avaliação de Clustering Pipeline}: Implemente um pipeline completo: (1) normalizar dados, (2) reduzir com PCA para 2D, (3) aplicar K-Means e HDBSCAN, (4) visualizar com t-SNE, (5) calcular todas as métricas de avaliação disponíveis. Reporte e discuta os resultados para o dataset Fashion-MNIST (10.000 amostras).
\end{enumerate}

% ============================================================
\section{Referências}
% ============================================================

\begin{thebibliography}{99}

\bibitem{agrawal1994}
Agrawal, R.; Srikant, R.
\newblock Fast algorithms for mining association rules.
\newblock \textit{Proceedings of the 20th VLDB Conference}, pp.~487--499, 1994.

\bibitem{han2000}
Han, J.; Pei, J.; Yin, Y.
\newblock Mining frequent patterns without candidate generation.
\newblock \textit{ACM SIGMOD Record}, 29(2):1--12, 2000.

\bibitem{macqueen1967}
MacQueen, J.
\newblock Some methods for classification and analysis of multivariate observations.
\newblock \textit{Proceedings of the 5th Berkeley Symposium on Mathematical Statistics and Probability}, 1:281--297, 1967.

\bibitem{arthur2007}
Arthur, D.; Vassilvitskii, S.
\newblock k-means++: The advantages of careful seeding.
\newblock \textit{Proceedings of the 18th Annual ACM-SIAM Symposium on Discrete Algorithms (SODA)}, pp.~1027--1035, 2007.

\bibitem{ester1996}
Ester, M.; Kriegel, H.-P.; Sander, J.; Xu, X.
\newblock A density-based algorithm for discovering clusters in large spatial databases with noise.
\newblock \textit{Proceedings of the 2nd ACM International Conference on Knowledge Discovery and Data Mining (KDD)}, pp.~226--231, 1996.

\bibitem{mcinnes2017}
McInnes, L.; Healy, J.; Astels, S.
\newblock hdbscan: Hierarchical density based clustering.
\newblock \textit{Journal of Open Source Software}, 2(11):205, 2017.

\bibitem{breunig2000}
Breunig, M. M.; Kriegel, H.-P.; Ng, R. T.; Sander, J.
\newblock LOF: Identifying density-based local outliers.
\newblock \textit{ACM SIGMOD Record}, 29(2):93--104, 2000.

\bibitem{liu2008}
Liu, F. T.; Ting, K. M.; Zhou, Z.-H.
\newblock Isolation forest.
\newblock \textit{IEEE International Conference on Data Mining (ICDM)}, pp.~413--422, 2008.

\bibitem{rousseeuw1987}
Rousseeuw, P. J.
\newblock Silhouettes: A graphical aid to the interpretation and validation of cluster analysis.
\newblock \textit{Journal of Computational and Applied Mathematics}, 20:53--65, 1987.

\bibitem{davies1979}
Davies, D. L.; Bouldin, D. W.
\newblock A cluster separation measure.
\newblock \textit{IEEE Transactions on Pattern Analysis and Machine Intelligence}, 1(2):224--227, 1979.

\bibitem{dempster1977}
Dempster, A. P.; Laird, N. M.; Rubin, D. B.
\newblock Maximum likelihood from incomplete data via the EM algorithm.
\newblock \textit{Journal of the Royal Statistical Society. Series B}, 39(1):1--38, 1977.

\bibitem{ng2002}
Ng, A. Y.; Jordan, M. I.; Weiss, Y.
\newblock On spectral clustering: Analysis and an algorithm.
\newblock \textit{Advances in Neural Information Processing Systems}, 14, 2002.

\bibitem{zhao2019}
Zhao, Y.; Nasrullah, Z.; Li, Z.
\newblock PyOD: A Python toolbox for scalable outlier detection.
\newblock \textit{Journal of Machine Learning Research}, 20(96):1--7, 2019.

\bibitem{caron2021}
Caron, M.; Touvron, H.; Misra, I.; et al.
\newblock Emerging properties in self-supervised vision transformers.
\newblock \textit{International Conference on Computer Vision (ICCV)}, 2021.

\end{thebibliography}

\end{document}
